\documentclass[11pt,a4paper]{article}

\usepackage[polish]{babel}
\usepackage[utf8]{inputenc}
\usepackage{polski}
\usepackage[T1]{fontenc}
\usepackage{indentfirst}
\usepackage{wrapfig}    % for wrapping figures, tables
\usepackage{isotope}
\usepackage{hyperref}
\usepackage{amsmath}
\usepackage{amssymb}
\usepackage{bm}
\usepackage{gensymb}
\usepackage{epsfig}
\usepackage{graphics}
\usepackage[shortlabels]{enumitem}
\usepackage{xspace}
\xspaceaddexceptions{[]\{\}}

%\def\hidesolutions{}  %← show solutions. Comment out to show solutions

%
%
%fixpagesize
\pagestyle{empty}
\addtolength{\textwidth}{4cm}
\addtolength{\textheight}{6cm}
\addtolength{\evensidemargin}{-3cm}
\addtolength{\oddsidemargin}{-2cm}
\addtolength{\topmargin}{-3cm}
\parindent=0cm

%
%
% Changes figure placing algorithm
\renewcommand{\topfraction}{1}       % maximal fraction of a page allowed for figures
\renewcommand{\textfraction}{0.15}   % minimal number of text for figure-text shared pages
\renewcommand{\floatpagefraction}{0.95} % if two above does not help, this could do the job 
                                        % must be: floatpagefraction < topfraction !!!!
\renewcommand{\textfraction}{0} % minimum fraction of page, which must be  devoted to text
\renewcommand{\topfraction}{1}  % maximum fraction at top, which can be occupied whit floats
\setcounter{totalnumber}{400}   % increase the number of floats for one page
\setcounter{topnumber}{200}     % at all/top/bottom.
\setcounter{bottomnumber}{200}  %

%
%
%small distance in list/item/enum for enumitem package
\setlist[itemize,enumerate]{topsep=0em}
\setlist{noitemsep}

%Nuclear notations: usage \nucl{235}{92}{U}. Math mode optional
\newcommand{\nucl}[3]{\ensuremath{
  \phantom{\ensuremath{^{#1}_{#2}}}
  \llap{\ensuremath{^{#1}}}
  \llap{\ensuremath{_{\rule{0pt}{.75em}#2}}}
  \mbox{#3} } 
} 

\newcounter{zadanie}\newcommand{\zadanie}[1][]{\addtocounter{zadanie}{1} ~\\  {\bf \emph{Exercise \arabic{zadanie} #1 }} \\}
\newcounter{zaddom}\newcommand{\zaddom}[1][]{\addtocounter{zaddom}{1} ~\\  {\bf \emph{Homework \arabic{zaddom} #1 }} \\}

%dynamic solutions hiding
\usepackage{comment}

\newif\ifhidesolutions
\ifdefined\hidesolutions
  \hidesolutionstrue
\else
  \hidesolutionsfalse
\fi

\ifhidesolutions
  \excludecomment{solution}
\else
  \newenvironment{solution}{\vspace{1cm} \par\textbf{Solution.} \vspace{1cm}}{\newpage}
\fi
%%%%%%%%%%%%%%%%%%%%%%%%%%%%%%%%%%%%%%%%%%
\begin{document}         

%%%%%%%%%%%%%%%%%%%%%%%%%%%%%%%%%%%%%%%%%%%%%%%%%%%%%
\begin{centering}
  {\bf {\Large
 Introduction to Experimental Particle Physics I \\
  Set XI: accelerators \\
  }}   
\end{centering}

\vspace{1cm}
%%%%%%%%%%%%%%%%%%%%%%%%%%%%%%%%%%%%%%%%%%%%%%%%%%%%%%%
\zadanie[Synchrotron radiation]

Find the rletivistic (=involving $dp^{\mu}/d\tau$) Larmor formula for power radiated by a charge under arbitrary acceleration. 
Derive formulas for two limiting cases:
\begin{itemize}
  \item $P_{L}$ - longitudinal acceleration,
  \item $P_{T}$ - transverse acceleration, with negligible energy gain
\end{itemize}

\vspace{0.5cm}

{\bf Hint:} assume particle is relativistic: 
$\beta \simeq 1$, 

\vspace{0.5cm}

Then calculate energy loss per turn in MeV for:
\begin{enumerate}
  \item proton in LHC at the top energy of 6.3 TeV
  \item electron in LEP at the top energy of 104 GeV
\end{enumerate}

Assume:
\begin{itemize}
  \item bending radius $\rho = 2803$ m 
  \item accelerating gradient $dE/dx = 5$ MeV/m
\end{itemize}

{\bf Hint:} create Python notebook for calculations. 


\begin{solution}

  Larmor formula for power radiated by a charge under acceleration:
\begin{align*}
  P = -\frac{e^{2}}{4\pi \epsilon_0} \frac{2}{3} \frac{1}{m^{2}c^3}  \frac{dp^{\mu}}{d\tau}\frac{dp_{\mu}}{d\tau} 
\end{align*}

Linear acceleration:
\begin{align*}
  & \frac{dp^{\mu}}{d\tau}\frac{dp_{\mu}}{d\tau} = \left[\frac{1}{c^{2}} \frac{dE}{d\tau}\frac{dE}{d\tau} - \frac{d\vec{p}}{d\tau}\cdot\frac{d\vec{p}}{d\tau}\right] = \\
  & || \tau = \gamma t \qquad |\vec{p}| = \beta E \\
  &    m^{2} = (E/c)^{2} - p^{2} \to E/c^{2}dE = pdp\\
  & dp = \frac{E}{p c^{2}} dE =  \frac{1}{\beta c^{2}} dE || \\
  & \frac{\gamma^{2}}{c^{2}} \left[\left( \frac{dE}{dt} \right)^{2} - \frac{1}{\beta^{2}} \left(\frac{dE}{dt} \right)^{2} \right] = \\
  & \frac{1}{c^{2}} \frac{\gamma^{2}}{\beta^{2}} \left( \frac{dE}{dt} \right)^{2} (\beta^{2} - 1) = -\frac{1}{\beta^{2}c^{2}} \left( \frac{dE}{dt} \right)^{2}
\end{align*}

So:
\begin{align*}
  & P_{L} = \frac{e^{2}}{4\pi \epsilon_0} \frac{2}{3} \frac{1}{m^{2}\beta^{2}c^5}  \left( \frac{dE}{dt} \right)^{2} = \\
  & || \beta \simeq 1 \to  ct = x|| \\
  & P_{L} = \frac{e^{2}}{4\pi \epsilon_0} \frac{2}{3} \frac{1}{m^{2}c^3}  \left( \frac{dE}{dx} \right)^{2}  
\end{align*}


Transverse acceleration:
\begin{align*}
  & \frac{dp^{\mu}}{d\tau}\frac{dp_{\mu}}{d\tau} = \frac{dE}{d\tau}\frac{dE}{d\tau} - \frac{d\vec{p}}{d\tau}\cdot\frac{d\vec{p}}{d\tau} = \\
  & || \tau = \gamma t \qquad dE \ll dp \qquad \frac{d\vec{p}}{dt} = \vec{\omega} \times \vec{p} || = \\
 & - \gamma^{2} (\omega \cdot |\vec{p}|)^{2} 
\end{align*}

So: 
\begin{align*}
  & P_{T} = \frac{e^{2}}{4\pi \epsilon_0} \frac{2}{3} \frac{\gamma^{2}}{m^{2}c^3}  (\omega \cdot |\vec{p}|)^{2} = \\
  & || \omega = \frac{\beta c}{\rho} \qquad |\vec{p}| = \gamma \beta c m|| \\
  & P_{T} = \frac{e^{2}}{4\pi \epsilon_0} \frac{2}{3} \frac{\gamma^{4} \beta^{4}c}{\rho^{2}}
\end{align*}

Energy lost per turn:
\begin{align*}
  & E_{L} = P_{L} \cdot \frac{2\pi \rho}{\beta c} = \frac{e^{2}}{3\epsilon_0} \frac{\rho}{m^{2}c^{4} \beta}  \left( \frac{dE}{dx} \right)^{2} \\
  & E_{T} = P_{T} \cdot \frac{2\pi \rho}{\beta c} = \frac{e^{2}}{3\epsilon_0} \frac{\gamma^{4} \beta^{3}}{\rho}
\end{align*}


Numerical calculations for LHC and LEP:
\begin{align*}
  & E_{L}^{LHC} = \frac{e^{2}}{3\epsilon_0} \frac{\rho}{m_{p}^{2}c^{4} \beta}  \left( \frac{dE}{dx} \right)^{2} \simeq 4\cdot 10^{-16} \text{ MeV} \\
  & E_{T}^{LHC} = \frac{e^{2}}{3\epsilon_0} \frac{\gamma^{4} \beta^{3}}{\rho} \simeq 5 \cdot 10^{-3} \text{ MeV} \\
  & E_{L}^{LEP} = \frac{e^{2}}{3\epsilon_0} \frac{\rho}{m_{e}^{2}c^{4} \beta}  \left( \frac{dE}{dx} \right)^{2} \simeq 1.6 \cdot 10^{-9} \text{ MeV} \\
  & E_{T}^{LEP} = \frac{e^{2}}{3\epsilon_0} \frac{\gamma^{4} \beta^{3}}{\rho} \simeq 4956 \text{ MeV} \\
\end{align*}

\end{solution}
%%%%%%%%%%%%%%%%%%%%%%%%%%%%%%%%%%%%%%%%%%%%%%%%%%%%%%%
%%%%%%%%%%%%%%%%%%%%%%%%%%%%%%%%%%%%%%%%%%%%%%%%%%%%%%%
\zadanie[M\o ller factor]

The cross section is defined as ratio of the reaction 
rate to the flux, F, of incoming particles:
\begin{align*}
  \sigma = \frac{R}{F}
\end{align*}

Reaction rate per volume and time element is:
\begin{align*}
  R = \frac{dN}{dV dt} = \frac{dN}{d^{3}x dt} = \frac{dN}{d^{4}x}
\end{align*}

The cross section is a Lorentz invariant quantity. Please:
\begin{itemize}
  \item show that the reaction rate is a Lorentz invariant quantity,
  \item read about the M\o ller factor, and a Lorentz invariant flux
\end{itemize}




\begin{solution}

{\bf Reaction rate:} the numerator is just a number of occurrences, which is Lorentz invariant. The denominator is a volume element in 4D space-time, which is also Lorentz invariant. Therefore, the reaction rate is a Lorentz invariant quantity.
Denominator is a volume element in 4D space-time, which is also Lorentz invariant, 
as can be shown using the Jacobian of the Lorentz transformation: 
\begin{align*}
  d^{4}x' = \det(\Lambda) d^{4}x = | \lambda \in SO(1, 3)  \to det(\Lambda)=1 | = d^{4}x
\end{align*}

Therefore, the reaction rate is a Lorentz invariant quantity.

\vspace{1cm}

{\bf The flux:}
Nonrelativistc flux is defined as:
\begin{align*}
  F = n_{1} n_{2} v_{12}
\end{align*}
where $n_{1}$ and $n_{2}$ are the number densities of the two colliding particle beams, 
and $v_{12}$ is the {\bf observed} velocity which two beams approach each one in the laboratory frame. 
This is not relative velocity, and can be greater than $c$. Let us construct the Lorentz invariant flux that will be equal to nonrelativistic flux 
in the LAB frame.

We start from the target rest frame, denoted with star, where 
we have $v_{rel} = v_{1}^{*}$, and $n_{1}^{*}, n_{2}^{*}$ are densities of the two colliding particle beams in the target rest frame.
The value of the flux in the target rest frame is:
\begin{align*}
  F = n_{1}^{*} n_{2}^{*} v_{1}^{*}
\end{align*}

We want $F$ to be Lorentz invariant, so we have to express it by invariant quantities.
Particle density transformation from rest frame to LAB frame is:
\begin{align*}
  & n^{0} = \frac{N}{V} = \frac{N}{A \cdot L} \\
  & n^{LAB} = \frac{N}{A \cdot L^{LAB}} = \frac{N}{A \cdot L} \frac{L}{L^{LAB}} = \\
  & || L^{LAB} = L \gamma || \to \\
  & n^{LAB} = \gamma n^{0}
\end{align*}

Beam densities in the target rest frame:
\begin{align*}
  & n_{1}^{*} = \frac{n_{1}^{LAB}}{\gamma_{1}} \\
  & n_{2}^{*} = n^{0} \gamma_{12} = \frac{n_{2}^{LAB}}{\gamma_{2}} \gamma_{12} 
\end{align*}

Relative velocity in the target rest frame is:
\begin{align*}
  v_{rel} = v_{1}^{*} = \frac{\sqrt{(\vec{v}_{1} - \vec{v}_{2})^{2} - (\vec{v}_{1} \times \vec{v}_{2})^2}}{1 - \vec{v}_{1}\vec{v}_{2}}
\end{align*}

and corresponding Lorentz factor is:
\begin{align*}
  \gamma_{rel} = \frac{1}{\sqrt{1 - v_{rel}^{2}}} = 
  \gamma_{1} \gamma_{2} (1 - \vec{v}_{1} \cdot \vec{v}_{2})
\end{align*}

Finally, the flux in the target rest frame is:
\begin{align*}
  F = \frac{n_{1}^{LAB}}{\gamma_{1}}\frac{n_{2}^{LAB}}{\gamma_{2}} \gamma_{1} \gamma_{2} (1 - \vec{v}_{1} \cdot \vec{v}_{2}) v_{rel} = \\
  = n_{1}^{LAB} n_{2}^{LAB}\sqrt{(\vec{v}_{1} - \vec{v}_{2})^{2} - (\vec{v}_{1} \times \vec{v}_{2})^2}
\end{align*}

\vspace{0.2cm}

{\bf Note:} the relativistic flux is equal to the nonrelativistic flux in the LAB frame
for the case of collinear beams, where $\vec{v}_{1} \times \vec{v}_{2} = 0$.

\vspace{0.5cm}

{\bf References:}
\href{https://cds.cern.ch/record/941318/files/p361.pdf}{Concept of liminosity},~
\href{https://www.osti.gov/servlets/purl/836235}{The M\o ller luminosity factor}

\end{solution}

%%%%%%%%%%%%%%%%%%%%%%%%%%%%%%%%%%%%%%%%%%%%%%%%%%%%%%%
%%%%%%%%%%%%%%%%%%%%%%%%%%%%%%%%%%%%%%%%%%%%%%%%%%%%%%%
\end{document}
